\begin{proposition}[谱四次曲线 Jacobian 的几何单纯性：不可约 Frobenius 多项式证书]\label{prop:fold-gauge-anomaly-spectral-quartic-jacobian-simple}
\leanverified{paper\_fold\_gauge\_anomaly\_spectral\_quartic\_jacobian\_simple}
沿用命题 \ref{prop:fold-gauge-anomaly-spectral-quartic-geometry} 的光滑平面四次曲线 \(C/\QQ\)，并令 \(J:=\operatorname{Jac}(C)\)。
则 \(J\) 在 \(\overline{\QQ}\) 上单纯。
更具体地，在素数 \(p=7\) 处 \(C\) 具有良好约化，且其约化曲线 \(C/\FF_{7}\) 的局部 \(L\)-多项式满足
\[
\boxed{\ L_{7}(T)=1+4T+9T^{2}+9T^{3}+63T^{4}+196T^{5}+343T^{6}\ }
\]
并在 \(\QQ[T]\) 中不可约。
\end{proposition}
\begin{proof}[证明要点]
脚本 \texttt{scripts/exp\_fold\_gauge\_anomaly\_spectral\_quartic\_jacobian\_L7\_audit.py} 以有限域穷举实现如下核验：
\begin{enumerate}
  \item 在 \(p=7\) 上对齐次模型 \(F^{\mathrm{hom}}(\mu,u,w)=0\) 检验无奇点，从而 \(C\) 在 \(p=7\) 处良好约化。
  \item 记 \(N_n:=\#C(\FF_{7^{n}})\ (n=1,2,3)\)，并由点数恒等式
  \(
  N_n=7^n+1-\sum_{i=1}^{6}\alpha_i^{n}
  \)
  回收幂和 \(S_n:=\sum_{i=1}^{6}\alpha_i^{n}\)；
  再由 Newton 恒等式确定 \(L_7(T)=\prod_{i=1}^{6}(1-\alpha_iT)\) 的系数，并得到上述闭式。
  \item 将 \(L_7(T)\) 模 \(13\) 化后在 \(\FF_{13}[T]\) 中保持不可约，从而 \(L_7(T)\) 在 \(\QQ[T]\) 中不可约。
\end{enumerate}
相应的可引用 TeX 证书片段为：
% Auto-generated by scripts/exp_fold_gauge_anomaly_spectral_quartic_jacobian_L7_audit.py
\[
\#C(\FF_{7})=12,\qquad \#C(\FF_{49})=52,\qquad \#C(\FF_{343})=327.
\]
\[
L_{7}(T)=343 T^{6} + 196 T^{5} + 63 T^{4} + 9 T^{3} + 9 T^{2} + 4 T + 1.
\]
\[
L_{7}(T)\bmod 13\ \text{在}\ \FF_{13}[T]\ \text{中不可约}\qquad(\text{因子次数分解：} [6]).
\]

\par
由 Honda--Tate 理论，\(L_7(T)\) 的不可约性蕴含约化雅可比簇 \(J_{\FF_7}\) 在 \(\overline{\FF}_7\) 上为简单阿贝尔三维簇，故其端同态代数无非平凡幂等元。
另一方面，若 \(J_{\overline{\QQ}}\) 可分，则 \(\mathrm{End}^0_{\overline{\QQ}}(J)\) 含非平凡幂等元；由良好约化处端同态专门化映射的单射性（Tate--Faltings 型结果），该幂等元将注入到 \(\mathrm{End}^0_{\overline{\FF}_{7}}(J_{\FF_7})\)，矛盾。
故 \(J\) 在 \(\overline{\QQ}\) 上单纯。证毕。
\end{proof}

\begin{proposition}[第二良素处的局部 \texorpdfstring{$L$}{L}-多项式]\label{prop:fold-gauge-anomaly-spectral-quartic-jacobian-L13}
\leanverified{paper\_fold\_gauge\_anomaly\_spectral\_quartic\_jacobian\_l13}
在素数 \(p=13\) 处，曲线 \(C\) 亦具有良好约化，且其约化曲线 \(C/\FF_{13}\) 的局部 \(L\)-多项式满足
\[
\boxed{\ L_{13}(T)=1+3T+7T^{2}+4T^{3}+91T^{4}+507T^{5}+2197T^{6}\ }
\]
并在 \(\QQ[T]\) 中不可约；其三次系数 \(4\) 不被 \(13\) 整除，从而约化雅可比簇 \(J_{\FF_{13}}\) 为普通阿贝尔三维簇。
\end{proposition}
\begin{proof}[证明要点]
脚本 \texttt{scripts/exp\_fold\_gauge\_anomaly\_spectral\_quartic\_jacobian\_L13\_audit.py} 以有限域穷举核验良好约化与点计数，并由 Newton 恒等式回收 \(L_{13}(T)\) 的闭式；
其在 \(\QQ[T]\) 中的不可约性由模 \(43\) 的不可约性证书给出。
相应可引用的 TeX 片段为：
% Auto-generated by scripts/exp_fold_gauge_anomaly_spectral_quartic_jacobian_L13_audit.py
\[
\#C(\FF_{13})=17,\qquad \#C(\FF_{169})=175,\qquad \#C(\FF_{2197})=2174.
\]
\[
L_{13}(T)=2197 T^{6} + 507 T^{5} + 91 T^{4} + 4 T^{3} + 7 T^{2} + 3 T + 1.
\]
\[
L_{13}(T)\bmod 43\ \text{在}\ \FF_{43}[T]\ \text{中不可约}\qquad(\text{因子次数分解：} [6]).
\]

证毕。
\end{proof}

\begin{theorem}[几何端同态环的极小性与绝对单纯性]\label{thm:fold-gauge-anomaly-spectral-quartic-jacobian-endomorphism-Z}
\leanverified{paper\_fold\_gauge\_anomaly\_spectral\_quartic\_jacobian\_endomorphism\_z}
在上述记号下，
\[
\operatorname{End}_{\overline{\QQ}}(J)\cong \ZZ,
\qquad
\text{因而 }J\text{ 在 }\overline{\QQ}\text{ 上绝对单纯。}
\]
\end{theorem}
\begin{proof}
取 \(p\in\{7,13\}\)。
由命题 \ref{prop:fold-gauge-anomaly-spectral-quartic-jacobian-simple} 与命题 \ref{prop:fold-gauge-anomaly-spectral-quartic-jacobian-L13}，\(C\) 在 \(p\) 处良好约化且 \(L_p(T)\) 在 \(\QQ[T]\) 中不可约，并且 \(T^{3}\) 系数不被 \(p\) 整除。
由 Honda--Tate 理论，约化雅可比簇 \(J_{\FF_p}\) 在 \(\overline{\FF}_p\) 上为普通且简单的阿贝尔三维簇，并满足
\[
\operatorname{End}^0_{\overline{\FF}_p}(J_{\FF_p})\cong K_p,
\qquad
K_p:=\QQ(\pi_p),\qquad [K_p:\QQ]=6,
\]
其中 \(\pi_p\) 为 Frobenius 元。
\par
另一方面，良好约化处端同态的专门化给出单射
\[
\operatorname{End}^0_{\overline{\QQ}}(J)\hookrightarrow \operatorname{End}^0_{\overline{\FF}_p}(J_{\FF_p})\cong K_p,
\]
从而
\[
\operatorname{End}^0_{\overline{\QQ}}(J)\subseteq K_7\cap K_{13}.
\]
\par
记
\[
f_7(x):=x^6+4x^5+9x^4+9x^3+63x^2+196x+343,
\qquad
f_{13}(x):=x^6+3x^5+7x^4+4x^3+91x^2+507x+2197,
\]
它们分别为 \(L_7,L_{13}\) 的互反多项式，生成相应的 CM 域 \(K_7,K_{13}\)。
直接计算得
\[
\gcd\!\bigl(\mathrm{Disc}(f_7),\mathrm{Disc}(f_{13})\bigr)=9.
\]
其中 \(\mathrm{Disc}(f_p)\) 为整系数多项式判别式，等于整环阶 \(\ZZ[\pi_p]\) 的判别式，因而被 \(\mathrm{disc}(K_p)\) 所整除。
若 \(K_7\cap K_{13}\neq\QQ\)，则其非平凡子域的域判别式同时整除 \(\mathrm{disc}(K_7)\) 与 \(\mathrm{disc}(K_{13})\)，从而亦整除 \(9\)，故只能为 \(\QQ(\sqrt{-3})\)。
\par
但对素数 \(\ell=41\equiv 2\pmod 3\)，多项式 \(f_7\bmod 41\) 在 \(\FF_{41}[x]\) 中分解型为 \((3,3)\)，从而 \(41\) 在 \(K_7\) 中存在残基次数为 \(3\) 的素因子；
若 \(\QQ(\sqrt{-3})\subseteq K_7\)，则 \(\ell\equiv 2\pmod 3\) 在 \(\QQ(\sqrt{-3})\) 中惰性，故 \(K_7\) 上一切残基次数应为偶数，矛盾。
同理，对 \(\ell=11\equiv 2\pmod 3\)，有 \(f_{13}\bmod 11\) 的分解型为 \((1,1,2,2)\)，从而 \(11\) 在 \(K_{13}\) 中存在残基次数为 \(1\) 的素因子，亦排除 \(\QQ(\sqrt{-3})\subseteq K_{13}\)。
故必有 \(K_7\cap K_{13}=\QQ\)。
\par
于是 \(\operatorname{End}^0_{\overline{\QQ}}(J)\subseteq\QQ\)，从而 \(\operatorname{End}^0_{\overline{\QQ}}(J)=\QQ\) 且 \(\operatorname{End}_{\overline{\QQ}}(J)=\ZZ\)。
最后，若 \(J\) 在 \(\overline{\QQ}\) 上可分，则 \(\operatorname{End}^0_{\overline{\QQ}}(J)\) 含非平凡幂等元，与其等于 \(\QQ\) 矛盾。
故 \(J\) 绝对单纯。证毕。
\end{proof}

\begin{corollary}[低属靶曲线的态射排除与 N\'eron--Severi 秩]\label{cor:fold-gauge-anomaly-spectral-quartic-no-low-genus-maps-ns-rank1}
对任意 \(\overline{\QQ}\)-定义的属 \(1\) 或属 \(2\) 曲线 \(D\)，不存在非恒定态射 \(C\to D\)。
并且 \(\mathrm{NS}(J)\) 的秩为 \(1\)。
\end{corollary}
\leanverified{paper\_fold\_gauge\_anomaly\_spectral\_quartic\_no\_low\_genus\_maps\_ns\_rank1}
\begin{proof}
若存在非恒定态射 \(C\to D\) 且 \(g(D)\in\{1,2\}\)，则诱导非零同态 \(J\to \operatorname{Jac}(D)\)，从而 \(J\) 存在维数 \(1\) 或 \(2\) 的非平凡阿贝尔商，进而在 \(\overline{\QQ}\) 上非单纯，矛盾于定理 \ref{thm:fold-gauge-anomaly-spectral-quartic-jacobian-endomorphism-Z}。
故第一断言成立。
\par
第二断言由标准同一 \(\mathrm{NS}(J)\otimes\QQ\cong \operatorname{End}^0_{\overline{\QQ}}(J)^{\dagger}\)（Rosati 对称端同态）与
\(\operatorname{End}^0_{\overline{\QQ}}(J)=\QQ\) 直接推出。证毕。
\end{proof}

\begin{conjecture}[最大 Mumford--Tate 与 \texorpdfstring{$\ell$}{l}-进满像]\label{conj:fold-gauge-anomaly-spectral-quartic-mumford-tate-gsp6}
设 \(\rho_{\ell}:\mathrm{Gal}(\overline{\QQ}/\QQ)\to \mathrm{GSp}_{6}(\ZZ_\ell)\) 为 \(J\) 的 \(\ell\)-进 Tate 模表示。
预期存在有限集合 \(S\) 使对一切 \(\ell\notin S\)，\(\rho_\ell\) 的像为有限指数子群，且其 Zariski 闭包为 \(\mathrm{GSp}_{6}\)。
等价地，\(J\) 的 Mumford--Tate 群为 \(\mathrm{GSp}_6\)，并呈 \(\mathrm{USp}(6)\) 型的 Sato--Tate 行为。
\end{conjecture}

\endinput
